\documentclass[uplatex,dvipdfmx]{jlreq}
\usepackage[fleqn,tbtags]{mathtools}
\usepackage{jlreq-deluxe}
\usepackage[noalphabet]{pxchfon} % must be after otf package
\usepackage{stix2} %欧文＆数式フォント
\usepackage[fleqn,tbtags]{mathtools} % 数式関連 (w/ amsmath)
\usepackage{url}
\begin{document}
\title{事前課題提出}
\author{25G1110 早川和輝}
\date{\today}
\maketitle
\section{技術的・社会的背景}
\subsection{社会的な背景}
現代社会において，商業施設や公共空間における「空間演出」は，人々の滞在価値を高める重要な要素となっている~\cite{110}．
従来の音響システムは，あらかじめ録音された音源をスピーカーから一方向に流す受動的なものが主流であった．
しかし，近年の体験型アートの普及により，鑑賞者の動きや環境の変化をトリガーとして，リアルタイムに反応・変化するインタラクティブな音響空間が求められている~\cite{110}．
本プロジェクトでは，複数の音源を配置し，それらが「輪唱」という形で互いに追いかけ合う音響構造を構築する．
さらに，主旋律に加えてエンターテインメント性を高める装飾音を独立したデバイスから発生させることで，鑑賞者に対してエンターテイメント性の高い体験を提供することを目指す．

\subsection{技術的背景}
このような音響システムを実現するためには，arduinoを用いたリアルタイム制御と，デバイス間の低遅延な通信技術が不可欠である~\cite{mitubishi}．
本システムでは，1台の同期用親機と5台の再生用子機を無線で接続する「1対多」のスター型ネットワーク構成をとる．
この構成には以下の技術的利点がある．
第一に，配線の制約を排除することで，広範な空間内の自由な位置に音源を配置可能となる点である．
第二に，負荷分散の観点である．
センサーデータの取得と解析を同期用の親機に集約し，各子機は受信したトリガーに基づいて音楽再生に専念することで，arduinoの処理能力を最大限に引き出し，音楽再生の安定性を確保できる．
特に，本システムの核心である「輪唱」の実現においては，複数のマイコン間での厳密なタイミング制御が要求される．
無線通信を介して同期信号をパケット化し，各子機へ高速に配送する仕組みを構築することで，音楽的な調和を保った分散演奏が可能となる．
さらに，エンターテイメント性を高める装飾音の独立した再生は，鑑賞者の没入感を高めるための重要な要素である．
このためには，主旋律の演奏と装飾音の再生が干渉しないように，適切なタイミング制御と音量バランスの調整が必要である．

\section{既存のシステム}
従来の空間演出システムは，主に受動的な音響配置に依存しており，鑑賞者の動きや環境の変化に対するリアルタイムな反応が限定的であった~\cite{110}．
一方で，近年の体験型アートの普及に伴い，インタラクティブな音響空間の需要が高まっている．
きわめてシンプルな例としては，センサーを用いて鑑賞者の位置や動きを検知し，それに応じて音響を変化させるシステムがある．
しかし，これらのシステムは通常，単一の音源からの出力に依存しており，複数の音源が相互に追いかい合うような複雑な音響構造を実現することは困難であった．
さらに，装飾音の独立した再生に関しても，従来のシステムでは主旋律の演奏と装飾音の再生が干渉しやすく，鑑賞者の没入感を損なう可能性があった．
これに対して，本プロジェクトでは，複数の音源が「輪唱」という形で互いに追いかけ合う音響構造を構築し，さらに装飾音の独立した再生を実現することで，鑑賞者に対してエンターテイメント性の高い体験を提供することを目指す．
このためには，arduinoを用いたリアルタイム制御と，デバイス間の低遅延な通信技術が不可欠である~\cite{mitubishi}．
特に，複数のマイコン間での厳密なタイミング制御が要求される「輪唱」の実現においては，無線通信を介して同期信号をパケット化し，各子機へ高速に配送する仕組みを構築することで，音楽的な調和を保った分散演奏が可能となる．
さらに，装飾音の独立した再生に関しても，主旋律の演奏と装飾音の再生が干渉しないように，適切なタイミング制御と音量バランスの調整が必要である．

\section{提案するシステムの概要}
システムは，前述した1対多のワイヤレス通信ネットワークを基盤とし，鑑賞者の「拍手」という直感的な音響アクションをトリガーとして，聴覚と視覚の双方に訴えかけるインタラクティブな空間演出システムである．
親機となるarduinoには高感度なサウンドセンサーを搭載し，子機群には音声再生モジュールとLEDマトリクスディスプレイをそれぞれ割り当てる．
本システムの体験は，鑑賞者の拍手に応じて二つの段階で構成される．
第一段階として，空間が静寂な待機状態において鑑賞者が手を一回叩くと，親機のサウンドセンサーがその音圧のピークを検知する．
親機は環境ノイズと意図的な拍手音をプログラム上の閾値判定によって識別し，拍手であると認識した瞬間に特定の音声再生用子機に対して無線でトリガー信号を送信する．
これにより，空間内に最初の主旋律となる音楽の再生が開始される．
第二段階として，主旋律が再生されている最中に鑑賞者がもう一度手を叩くと，空間演出はより複雑でダイナミックなものへと移行する．
親機は二回目の拍手を検知すると，今度は待機していた複数の音声再生用子機に対して，それぞれ異なる遅延時間を設定した再生開始コマンドを一斉に送信する．
各子機は指定されたタイミングで順次音源の再生を開始するため，空間全体で音が互いに追いかけ合うような「輪唱」の構造が構築される．
さらに，この二回目の拍手による演出の拡張は聴覚だけにとどまらない．
親機は音声信号の送信と完全に同期して，LEDマトリクスを接続した視覚演出用の子機に対しても描画コマンドを送信する．
コマンドを受信した子機は，即座にマトリクス上にプログラムされた幾何学図形やパターンを描画する．
これにより，音楽が重層的な輪唱へと展開する劇的なタイミングに合わせて視覚的なアクセントが生み出される．
技術的には，これらの動作を破綻なく成立させるために親機側での厳密なステート（状態）管理が求められる．
親機は現在のシステムが「待機中」なのか「音楽再生中」なのかを常に記憶し，同じ拍手という入力に対しても，状態に応じて異なるパケットを生成して送信する分岐処理を行う．
このように，低遅延の無線通信と状態遷移のアルゴリズムを組み合わせることで，鑑賞者のシンプルな身体動作から，複雑な音響構造と視覚的変化が連動するエンターテインメント性の高い体験を実現する．
\begin{thebibliography}{9}
\bibitem{110}株式会社ワントゥーテン（1→10）．Location Based　Entertainment（LBE）．\url{https://www.1-10.com/business/lbe/}（参照2026年4月22日）
\bibitem{mitubishi}三菱電機株式会社．高信頼・低遅延な無線通信技術．青山・鈴木・梅田・木下・武. 特集論文. 29（405）.
\end{thebibliography}

\end{document} 
